\documentclass[12pt]{article}
\usepackage[utf8]{inputenc}
\usepackage[T1]{fontenc}
\usepackage{hyperref}
\title{Tamashek Language Resources for K-12 Curriculum Development}
\author{Compiled by Claude Sonnet 5 (claude-sonnet-5) for Andr\'as Kornai}
\date{}
\begin{document}
\maketitle

\section{Introduction}
ISO 639-3 \texttt{tmh} is the Tamashek (Tuareg) macrolanguage, a Berber language. The prompt gave only the code, so this report assumes that scope. The search results were mostly publisher catalogues and academic pages. They turned up no school textbooks or primers. Varieties differ by country (Mali, Niger, Burkina Faso, Algeria), so resources should be matched to the variety.

\section{Primary Educational Resources}
No native-language school primer or textbook turned up in the search results. Literacy programmes in the region are reported to use a Latin-based alphabet. I did not confirm which programmes.

\section{Linguistic and Reference Materials}
\begin{itemize}
\item Sudlow, \textit{Dictionary of the Tamasheq of North-East Burkina Faso}: \url{https://www.koeppe.de/titel_dictionary-of-the-tamasheq-of-north-east-burkina-faso}
\item Sudlow, \textit{The Tamasheq of North-East Burkina Faso: A Learner's Grammar}: \url{https://www.koeppe.de/titel_the-tamasheq-of-north-east-burkina-faso-a-learner-s-grammar-2}
\item Prasse, \textit{Tuareg Elementary Course (Tahaggart)}: \url{https://www.koeppe.de/titel_tuareg-elementary-course-tahaggart}
\item Webonary site for a Tamahaq dictionary: \url{https://www.webonary.org/tamahaq/overview/foreword/?lang=en}
\item Jeffrey Heath's Berber bibliography page (Tamashek of Mali dictionary, grammar and texts are listed there): \url{https://websites.umich.edu/~jheath/Berber.html}
\end{itemize}

\section{Religious and Community Resources}
None were found in the search results.

\section{Digital and Audio Resources}
The only digital resource found is the Webonary Tamahaq dictionary listed above. No audio resources were found.

\section{Curriculum Development Context}
The sources found are scholarly and learner-oriented. They are not classroom materials. A curriculum would need new, variety-specific materials, using the existing dictionaries and grammars as linguistic bases. Further searching for national literacy services and NGO programmes in Mali and Niger is recommended.
\end{document}